% GENERATED FILE. Edit canonical lessons, then run npm run generate:books.
% canonical-source-hash: fnv1a64:856790fee87e85b7
% canonical-lessons: MR-C419-vasti, MR-C419-vasahat, MR-C419-peth, MR-C419-bol, MR-C419-cal, MR-R419-review-412-419

\chapter{Neighbourhoods}
\label{ch:mr-neighbourhoods}
\hlchaptermodality{419}

\begin{chapteropening}
\textbf{By the end of this chapter:} \emph{I can talk about settlements, housing colonies, old quarters, lanes and tenements.}

Complete the last lesson of chapter 419: I can talk about settlements, housing colonies, old quarters, lanes and tenements.
\end{chapteropening}

\section[\texorpdfstring{\mr{वस्ती} (vastī)}{वस्ती (vastī)}]{\mr{वस्ती} (vastī) — a settlement; a neighbourhood}

\label{lesson:MR-C419-vasti}



\begin{quote}
Before the new one: say the Marathi for a small village, then the Marathi for a taluka.
\end{quote}

\subsection*{You'll want to know: \mr{वस्ती}}
\textbf{\mr{वस्ती}} — \emph{vastī} — \textquotedblleft{}a settlement; a neighbourhood\textquotedblright{}.

\begin{quote}
\mr{डोंगराखाली} \mr{एक} \mr{लहान} \mr{वस्ती} \mr{आहे}.
\end{quote}

\begin{quote}
\emph{ḍoṅgarākhālī ek lahān vastī āhe.}
\end{quote}

\begin{quote}
There is a small settlement at the bottom of the hill.
\end{quote}

Feminine.

Also the people who live there: \mr{सगळी} \mr{वस्ती} \mr{जमली}, the whole neighbourhood gathered.

\begin{cousinweb}
\textbf{\mr{वस्ती}} (\emph{vastī}) belongs with Marathi \emph{vasṇẽ}, \textquotedblleft{}to settle\textquotedblright{}, inherited from Sanskrit \emph{vasati}, \textquotedblleft{}stays, dwells\textquotedblright{} (CDIAL 11435); how the noun was formed is not settled.
\end{cousinweb}

\begin{grammarlens}[title={What you've built}]
One more word for getting around.
\end{grammarlens}

\subsection*{Your turn}
\begin{itemize}
\raggedright
  \item \emph{Say it:} \emph{vastī}
  \item \emph{Say it:} \emph{vastī}, once more
  \item \emph{Say it:} \emph{ḍoṅgarākhālī ek lahān vastī āhe.} — There is a small settlement at the bottom of the hill.
  \item \emph{Recall:} say the Marathi for a small village, then the Marathi for a taluka, then say \emph{vastī} again
\end{itemize}

\begin{tcolorbox}[breakable,colback=teal!4,colframe=teal!35!black,title={Before you move on}]
What does \mr{वस्ती} mean? (\textquotedblleft{}A settlement; a neighbourhood\textquotedblright{}.) Say it once more.
\end{tcolorbox}
\section[\texorpdfstring{\mr{वसाहत} (vasāhat)}{वसाहत (vasāhat)}]{\mr{वसाहत} (vasāhat) — a housing colony}

\label{lesson:MR-C419-vasahat}



\begin{quote}
Before the new one: say the Marathi for a taluka, then the Marathi for a settlement.
\end{quote}

\subsection*{You'll want to know: \mr{वसाहत}}
\textbf{\mr{वसाहत}} — \emph{vasāhat} — \textquotedblleft{}a housing colony\textquotedblright{}.

\begin{quote}
\mr{आमची} \mr{वसाहत} \mr{शहराच्या} \mr{बाहेर} \mr{आहे}.
\end{quote}

\begin{quote}
\emph{āmcī vasāhat śahrācyā bāher āhe.}
\end{quote}

\begin{quote}
Our housing colony is outside the city.
\end{quote}

Feminine.

Also a colony in history: a land settled and ruled from far away.

\begin{cousinweb}
\textbf{\mr{वसाहत}} (\emph{vasāhat}) is a Marathi word for a planned settlement; its origin is not settled.
\end{cousinweb}

\begin{grammarlens}[title={What you've built}]
One more word for getting around.
\end{grammarlens}

\subsection*{Your turn}
\begin{itemize}
\raggedright
  \item \emph{Say it:} \emph{vasāhat}
  \item \emph{Say it:} \emph{vasāhat}, once more
  \item \emph{Say it:} \emph{āmcī vasāhat śahrācyā bāher āhe.} — Our housing colony is outside the city.
  \item \emph{Recall:} say the Marathi for a taluka, then the Marathi for a settlement, then say \emph{vasāhat} again
\end{itemize}

\begin{tcolorbox}[breakable,colback=teal!4,colframe=teal!35!black,title={Before you move on}]
What does \mr{वसाहत} mean? (\textquotedblleft{}A housing colony\textquotedblright{}.) Say it once more.
\end{tcolorbox}
\section[\texorpdfstring{\mr{पेठ} (peṭh)}{पेठ (peṭh)}]{\mr{पेठ} (peṭh) — an old town quarter}

\label{lesson:MR-C419-peth}



\begin{quote}
Before the new one: say the Marathi for a settlement, then the Marathi for a housing colony.
\end{quote}

\subsection*{You'll want to know: \mr{पेठ}}
\textbf{\mr{पेठ}} — \emph{peṭh} — \textquotedblleft{}an old town quarter\textquotedblright{}.

\begin{quote}
\mr{पुण्यात} \mr{अनेक} \mr{जुन्या} \mr{पेठा} \mr{आहेत}.
\end{quote}

\begin{quote}
\emph{puṇyāt anek junyā peṭhā āhet.}
\end{quote}

\begin{quote}
Pune has many old quarters.
\end{quote}

Feminine. Plural: \mr{पेठा}, \emph{peṭhā}.

In Pune each old quarter is a \mr{पेठ}; the plural is \mr{पेठा}.

\begin{cousinweb}
\textbf{\mr{पेठ}} (\emph{peṭh}) is inherited: CDIAL lists Marathi \emph{peṭh}, \textquotedblleft{}town, quarter of a town, market\textquotedblright{}, under Sanskrit \emph{pratiṣṭhā}, \textquotedblleft{}foundation, homestead\textquotedblright{} (CDIAL 8608).
\end{cousinweb}

\begin{grammarlens}[title={What you've built}]
One more word for getting around.
\end{grammarlens}

\subsection*{Your turn}
\begin{itemize}
\raggedright
  \item \emph{Say it:} \emph{peṭh}
  \item \emph{Say it:} \emph{peṭh}, once more
  \item \emph{Say it:} \emph{puṇyāt anek junyā peṭhā āhet.} — Pune has many old quarters.
  \item \emph{Recall:} say the Marathi for a settlement, then the Marathi for a housing colony, then say \emph{peṭh} again
\end{itemize}

\begin{tcolorbox}[breakable,colback=teal!4,colframe=teal!35!black,title={Before you move on}]
What does \mr{पेठ} mean? (\textquotedblleft{}An old town quarter\textquotedblright{}.) Say it once more.
\end{tcolorbox}
\section[\texorpdfstring{\mr{बोळ} (boḷ)}{बोळ (boḷ)}]{\mr{बोळ} (boḷ) — a narrow lane}

\label{lesson:MR-C419-bol}



\begin{quote}
Before the new one: say the Marathi for a housing colony, then the Marathi for an old town quarter.
\end{quote}

\subsection*{You'll want to know: \mr{बोळ}}
\textbf{\mr{बोळ}} — \emph{boḷ} — \textquotedblleft{}a narrow lane\textquotedblright{}.

\begin{quote}
\mr{या} \mr{बोळातून} \mr{बाजारात} \mr{लवकर} \mr{पोहोचता} \mr{येतं}.
\end{quote}

\begin{quote}
\emph{yā boḷātūn bājārāt lavkar pohoctā yetaṁ.}
\end{quote}

\begin{quote}
You can reach the market quickly through this lane.
\end{quote}

Masculine.

Narrower than a \mr{गल्ली}, often a short cut.

\begin{cousinweb}
\textbf{\mr{बोळ}} (\emph{boḷ}) is a Marathi word for a narrow passage between houses; its origin is not settled.
\end{cousinweb}

\begin{grammarlens}[title={What you've built}]
One more word for getting around.
\end{grammarlens}

\subsection*{Your turn}
\begin{itemize}
\raggedright
  \item \emph{Say it:} \emph{boḷ}
  \item \emph{Say it:} \emph{boḷ}, once more
  \item \emph{Say it:} \emph{yā boḷātūn bājārāt lavkar pohoctā yetaṁ.} — You can reach the market quickly through this lane.
  \item \emph{Recall:} say the Marathi for a housing colony, then the Marathi for an old town quarter, then say \emph{boḷ} again
\end{itemize}

\begin{tcolorbox}[breakable,colback=teal!4,colframe=teal!35!black,title={Before you move on}]
What does \mr{बोळ} mean? (\textquotedblleft{}A narrow lane\textquotedblright{}.) Say it once more.
\end{tcolorbox}
\section[\texorpdfstring{\mr{चाळ} (cāḷ)}{चाळ (cāḷ)}]{\mr{चाळ} (cāḷ) — a tenement}

\label{lesson:MR-C419-cal}



\begin{quote}
Before the new one: say the Marathi for an old town quarter, then the Marathi for a narrow lane.
\end{quote}

\subsection*{You'll want to know: \mr{चाळ}}
\textbf{\mr{चाळ}} — \emph{cāḷ} — \textquotedblleft{}a tenement\textquotedblright{}.

\begin{quote}
\mr{शहरात} \mr{अजून} \mr{काही} \mr{जुन्या} \mr{चाळी} \mr{आहेत}.
\end{quote}

\begin{quote}
\emph{śahrāt ajūn kāhī junyā cāḷī āhet.}
\end{quote}

\begin{quote}
There are still some old tenements in the city.
\end{quote}

Feminine. Plural: \mr{चाळी}, \emph{cāḷī}.

A long building of one-room homes along a shared balcony; the plural is \mr{चाळी}.

\begin{cousinweb}
\textbf{\mr{चाळ}} (\emph{cāḷ}) is a Marathi word for a long tenement block; its origin is not settled.
\end{cousinweb}

\begin{grammarlens}[title={What you've built}]
One more word for getting around.
\end{grammarlens}

\subsection*{Your turn}
\begin{itemize}
\raggedright
  \item \emph{Say it:} \emph{cāḷ}
  \item \emph{Say it:} \emph{cāḷ}, once more
  \item \emph{Say it:} \emph{śahrāt ajūn kāhī junyā cāḷī āhet.} — There are still some old tenements in the city.
  \item \emph{Recall:} say the Marathi for an old town quarter, then the Marathi for a narrow lane, then say \emph{cāḷ} again
\end{itemize}

\begin{tcolorbox}[breakable,colback=teal!4,colframe=teal!35!black,title={Before you move on}]
What does \mr{चाळ} mean? (\textquotedblleft{}A tenement\textquotedblright{}.) Say it once more.
\end{tcolorbox}
\section[(dialogue)]{Looking back: from Finding the Way to Neighbourhoods}

\label{lesson:MR-R419-review-412-419}



\begin{quote}
Say the Marathi for a direction and a compass.
\end{quote}

\subsection*{Your turn}
Say these aloud:

\begin{itemize}
\raggedright
  \item a direction, a compass, a mark, to stray, to hang about
  \item lodging, a rest house, a guest house, an overnight stay, an eating house
  \item an honoured guest, hospitality, a reception desk, to ask after, to drop by
  \item an amenity, bedding, a woollen blanket, to spread out bedding, to cover oneself
  \item a peg, a washerman, to wring, a porter, to sprinkle
  \item a deposit, to settle down, a tenant, a strongbox, an elevator
  \item a town, a metropolis, a municipal council, a small village, a taluka
  \item a settlement, a housing colony, an old town quarter, a narrow lane, a tenement
\end{itemize}

\begin{tcolorbox}[breakable,colback=teal!4,colframe=teal!35!black,title={Before you move on}]
A direction? (\textbf{\mr{दिशा}}, \emph{diśā}.) A compass? (\textbf{\mr{होकायंत्र}}, \emph{hokāyantra}.) A mark? (\textbf{\mr{खूण}}, \emph{khūṇ}.) To stray? (\textbf{\mr{भरकटणे}}, \emph{bharkaṭṇe}.) To hang about? (\textbf{\mr{घुटमळणे}}, \emph{ghuṭmaḷṇe}.) Lodging? (\textbf{\mr{निवास}}, \emph{nivās}.) A rest house? (\textbf{\mr{धर्मशाळा}}, \emph{dharmaśāḷā}.) A guest house? (\textbf{\mr{विश्रामगृह}}, \emph{viśrāmgṛh}.) An overnight stay? (\textbf{\mr{मुक्काम}}, \emph{mukkām}.) An eating house? (\textbf{\mr{खानावळ}}, \emph{khānāvaḷ}.) An honoured guest? (\textbf{\mr{अतिथी}}, \emph{atithī}.) Hospitality? (\textbf{\mr{आदरातिथ्य}}, \emph{ādarātithya}.) A reception desk? (\textbf{\mr{स्वागतकक्ष}}, \emph{svāgatkakṣa}.) To ask after? (\textbf{\mr{विचारपूस} \mr{करणे}}, \emph{vicārpūs karṇe}.) To drop by? (\textbf{\mr{फिरकणे}}, \emph{phirakṇe}.) An amenity? (\textbf{\mr{सुविधा}}, \emph{suvidhā}.) Bedding? (\textbf{\mr{बिछाना}}, \emph{bichānā}.) A woollen blanket? (\textbf{\mr{घोंगडी}}, \emph{ghoṅgḍī}.) To spread out bedding? (\textbf{\mr{अंथरणे}}, \emph{antharṇe}.) To cover oneself? (\textbf{\mr{पांघरणे}}, \emph{pāṅgharṇe}.) A peg? (\textbf{\mr{खुंटी}}, \emph{khuṇṭī}.) A washerman? (\textbf{\mr{धोबी}}, \emph{dhobī}.) To wring? (\textbf{\mr{पिळणे}}, \emph{piḷṇe}.) A porter? (\textbf{\mr{हमाल}}, \emph{hamāl}.) To sprinkle? (\textbf{\mr{शिंपडणे}}, \emph{śimpaḍṇe}.) A deposit? (\textbf{\mr{अनामत}}, \emph{anāmat}.) To settle down? (\textbf{\mr{स्थायिक} \mr{होणे}}, \emph{sthāyik hoṇe}.) A tenant? (\textbf{\mr{भाडेकरू}}, \emph{bhāḍekarū}.) A strongbox? (\textbf{\mr{तिजोरी}}, \emph{tijorī}.) An elevator? (\textbf{\mr{उद्वाहक}}, \emph{udvāhak}.) A town? (\textbf{\mr{नगर}}, \emph{nagar}.) A metropolis? (\textbf{\mr{महानगर}}, \emph{mahānagar}.) A municipal council? (\textbf{\mr{नगरपालिका}}, \emph{nagarpālikā}.) A small village? (\textbf{\mr{खेडे}}, \emph{kheḍe}.) A taluka? (\textbf{\mr{तालुका}}, \emph{tālukā}.) A settlement? (\textbf{\mr{वस्ती}}, \emph{vastī}.) A housing colony? (\textbf{\mr{वसाहत}}, \emph{vasāhat}.) An old town quarter? (\textbf{\mr{पेठ}}, \emph{peṭh}.) A narrow lane? (\textbf{\mr{बोळ}}, \emph{boḷ}.) A tenement? (\textbf{\mr{चाळ}}, \emph{cāḷ}.)
\end{tcolorbox}
